% ======================================================================
% Parte 5: Diseñen Perfectópolis Segura
% ======================================================================

Como equipo, propongan 3 reglas fundamentales que deban aplicarse a todos los sistemas digitales de Perfectópolis, considerando temas como privacidad, gestión de contraseñas, cifrado, integridad, control de acceso, IA, monitoreo/respuesta o accesibilidad. Cada regla debe incluir su justificación y un ejemplo claro de aplicación.

\subsection*{Regla 1}
\begin{itemizeColor}
    \item \textbf{Regla:} 
    % TODO: Escribir la regla propuesta
    Todos los sistemas deben garantizar canales de acceso alternativos a los digitales. Tanto los sistemas digitales como los canales alternativos deben cumplir con estándares de accesibilidad (como el WCAG) para asegurar que los usuarios vulnerables a la tecnología tengan acceso a sus derechos.
    \item \textbf{Justificación:} 
    % TODO: Escribir la justificación
	    La digitalización que busca Perfectópolis puede marginar a usuarios sin acceso a los dispositivos para consumir sus plataformas, sin acceso a internet, con discapacidades o con problemas con el uso de la tecnología; en consecuencia, es posible que al digitalizar todo el sistema y no brindar alternativas para la población vulnerable, estos sean invisibilizados y no se pueda garantizar el ejercicio de sus derechos.

    \item \textbf{Ejemplo de aplicación:} 
    % TODO: Escribir el ejemplo de aplicación
	    El sistema de denuncias digital no debe reemplazar totalmente a un sistema de denuncias convencional, pues al carecer de otros medios para denunciar, a las poblaciones vulnerables se les hace más difícil, o imposible, denunciar. Por ello, el sistema de denuncias digital debe ser accesible aportando sistemas de accesibilidad al mismo tiempo que se proporcionan medios físicos como los ministerios públicos o sistemas de denuncia rápida conectados al sistema como los botones de pánico del C5 en México.
\end{itemizeColor}

\subsection*{Regla 2}
\begin{itemizeColor}
    \item \textbf{Regla:} 
    % TODO: Escribir la regla propuesta
Todos los sistemas deben brindar un cifrado apropiado de los datos personales y biométricos, tanto en el almacenamiento en bases de datos como en el tránsito entre servidores y clientes. Además, estos datos no se deben utilizar para otros fines ajenos a los especificados en un aviso de privacidad, ni ser almacenados por más tiempo del necesario.
    \item \textbf{Justificación:} 
    % TODO: Escribir la justificación
	    El almacenamiento inapropiado de datos, así como el que se almacenen por un tiempo muy prolongado, hace más probable su eventual filtración y hace más graves las consecuencias del robo de datos. Además, se requiere de transparencia en el uso de los datos para evitar que el propio gobierno, o empresas asociadas a los sistemas, puedan abusar de la información que obtienen de los sistemas.

    \item \textbf{Ejemplo de aplicación:} 
    % TODO: Escribir el ejemplo de aplicación
	    Los sistemas de vigilancia no pueden conservar las grabaciones por más de un tiempo determinado a menos que se demuestre que puedan ser útiles para la investigación de un delito; además, estas deben estar apropiadamente cifradas y protegidas ante cualquier alteración para evitar la modificación o eliminación de pruebas relevantes en un caso. De la misma manera, se debe ser transparente con el uso de los datos obtenidos de cualquiera de los sistemas de Perfectópolis y no utilizarse para fines distintos a los de los sistemas, como por ejemplo, el entrenamiento de modelos de inteligencia artificial o análisis de consumo.
\end{itemizeColor}

\subsection*{Regla 3}
\begin{itemizeColor}
    \item \textbf{Regla:} 
    % TODO: Escribir la regla propuesta
    Ningún sistema de inteligencia artificial puede tomar decisiones finales sin la revisión y aprobación exhaustiva de equipos de expertos humanos con una ética profesional que impida actuar sobre los propios intereses o los de una organización.
    \item \textbf{Justificación:} 
    % TODO: Escribir la justificación
	    Las herramientas de inteligencia artificial son aún herramientas relativamente nuevas y no han podido probar que carezcan de sesgos tecnológicos e ideológicos que puedan afectar la toma de decisiones; si bien es complicado probar que un humano no pueda tener estos mismos sesgos, al asegurar que se verifique con un equipo de expertos reducimos la probabilidad de que los sistemas obedezcan a intereses propios o de empresas en vez del bien común.

    \item \textbf{Ejemplo de aplicación:} 
    % TODO: Escribir el ejemplo de aplicación
	    La plataforma ParticipaPerfecto clasifica por medio de inteligencia artificial las denuncias; sin embargo, esta clasificación es evaluada por un equipo humano para verificar que es correcta. Además, se tiene un equipo humano con el que se pueda apelar la decisión para asegurar que no se invisibilicen denuncias que de manera sesgada la inteligencia artificial clasifique como poco importantes.
\end{itemizeColor}
